\begin{figure}[H]
\centering
\begin{tikzpicture}[
  x=1cm,
  y=1cm,
  figure block/.style={dissertation block,minimum height=1.55cm}
]
\node[figure block,text width=2.75cm] (src) at (-6,2) {Зовнішнє\\джерело};
\node[figure block,text width=2.90cm] (obj) at (-2,2) {Цифрове зображення\\або відеозапис};
\node[figure block,text width=3.25cm] (recv) at (2,2) {Приймання та\\форматна перевірка};
\node[figure block,text width=3.70cm] (proc) at (6,2) {Контрольоване розкодування\\та перетворення};
\node[figure block,text width=3.60cm] (orig) at (2,-0.8) {Оригінальний контейнер\\і службові дані};
\node[figure block,text width=3.25cm] (store) at (6,-0.8) {Зберігання та\\публікація};
\node[figure block,text width=4.20cm] (effect) at (6,-3.4) {Прояв\\прихованого впливу\\або формування ознак};

\coordinate (effectroute) at ($(proc.east)+(0.55,0)$);

\draw[dissertation arrow] (src.east) -- (obj.west);
\draw[dissertation arrow] (obj.east) -- (recv.west);
\draw[dissertation arrow] (recv.east) -- (proc.west);
\draw[dissertation arrow] (recv.south) -- (orig.north);
\draw[dissertation arrow] (orig.east) -- (store.west);
\draw[dissertation arrow] (proc.south) -- (store.north);
\draw[dissertation arrow] (store.south) -- (effect.north);
\draw[dissertation arrow] (proc.east) -- (effectroute)
  -- (effectroute |- effect.east) -- (effect.east);
\end{tikzpicture}
\caption{Маршрут цифрового зображення або відеозапису у вебсистемі та точки формування ознак прихованого впливу}
\label{fig:multimedia_threat_path}
\end{figure}
